\begin{lemma}[Strict shortening of a Rutgers counterexample]
Fix an original positive-rung extended Tashkinov tree \(\mathcal E\),
and assume \(\Delta(G)+1\leq k\).  Assume the following two assertions
uniformly for every maximum ETT \(\mathcal F\) having the same rung
number as \(\mathcal E\), with arbitrary certified input:
\begin{enumerate}
\item the displayed vertex set of \(\mathcal F\) is elementary under
  its generating coloring;
\item every proper coloring leaving exactly \(e_0\) uncolored and stable
  on the terminal ladder tree of \(\mathcal F\), relative to its
  terminal stored set and terminal series coloring, is a
  \(\noderef{ModColoring}\) of \(\mathcal F\).
\end{enumerate}
Let \(\mathcal P=(\sigma,U,\ldots)\) be a
\(\noderef{RutgersClosedPresentation}\) relative to \(\mathcal E\),
and let \(\mathcal C\) be a \(\noderef{RutgersCounterexample}\) on
\(\mathcal P\).  Write \(\mathcal R\) for the replacement ETT in
\(\mathcal P\), and \(T_n\) for its actual terminal ladder tree.  If
the middle ordered tree end of \(\mathcal C\) satisfies
\[
                              b_2\notin V(T_n),
\]
then there are another Rutgers closed presentation \(\mathcal P'\)
relative to the same fixed original \(\mathcal E\), and a Rutgers
counterexample \(\mathcal C'\) on \(\mathcal P'\), such that
\[
 \noderef{RutgersCounterexampleMeasure}(\mathcal C')
 <
 \noderef{RutgersCounterexampleMeasure}(\mathcal C).
\]
No minimality hypothesis on \(\mathcal C\) is assumed.
\end{lemma}
\begin{proof}
Put \(n=\mathcal R.\mathit{rounds}\), let \(T_n\) be its terminal
ladder tree, and let
\(D_n=\noderef{TashkinovStoredColors}(\mathcal R,n)\).  The dependent
series and final-coloring fields of
\(\noderef{CorrespondingExtendedTashkinovTree}\) identify \(\sigma\)
with the terminal series coloring of \(\mathcal R\).  Its
replacement-proper and replacement-one-uncolored fields therefore say
that \(\sigma\) is proper in the sense of
\(\noderef{IsProperEdgeColoring}\) and leaves exactly \(e_0\)
uncolored in the sense of \(\noderef{IsOneEdgeUncolored}\).

The terminal route in \(\mathcal P\) supplies both
\[
             T_n\preccurlyeq U
   \quad\hbox{and}\quad
             \noderef{ColoredVertexSetClosed}(G,\sigma,V(U)). \tag{1}
\]
Indeed, in the PE alternative the
\(\noderef{TashkinovClosure}\) starts at \(T_n\); in the SE/RE
alternative it starts at an exact one-edge extension of \(T_n\) supplied
by \(\noderef{TreeSequenceAdjoinEdge}\).  The augmentation sequence
inside either closure makes its source an exact initial segment of \(U\)
by \(\noderef{TashkinovAugmentationSequenceInitialSegment}\), and the
entrywise relations in \(\noderef{TreeSequenceInitialSegment}\) compose.
Applying \(\noderef{TreeSequenceInitialSegmentVertexSetSubset}\) gives
the corresponding vertex-set inclusions.  In particular all three
inside ends \(b_i\), which belong to \(V(U)\) by
\(\noderef{TreeExitPath}\), are genuine displayed vertices of \(U\).

By \(\noderef{MissingColorsNonemptyOfMaximumDegree}\), choose
\[
                    \gamma\in
                    \noderef{MissingColors}(G,\sigma,b_3).   \tag{2}
\]
Thus \(\gamma\in\noderef{TreeMissingColors}(\sigma,U)\).
The retained identity \(f_3\) has \(\sigma(f_3)=\beta\) by
\(\noderef{BoundaryColorEdges}\), and is incident with \(b_3\) by the
counterexample fields.  The witness definition
\(\noderef{PresentColors}\), together with (2), gives
\[
                              \gamma\ne\beta.                 \tag{3}
\]

Define
\[
 \mu_3=\noderef{OutsideTreeColorSwap}(U,\sigma,\alpha,\gamma).
                                                                    \tag{4}
\]
If \(\alpha=\gamma\), this is just \(\sigma\).  If
\(\alpha\ne\gamma\), both colors are tree-missing on \(U\): this is a
field of \(\mathcal C\) for \(\alpha\), and follows from (2) for
\(\gamma\).  Apply \(\noderef{OutsideTreeColorSwapStable}\) using the
closedness in (1).  In either equality branch, \(\mu_3\) is proper,
leaves exactly \(e_0\) uncolored, is
\((U,D_n,\sigma)\)-stable in the sense of
\(\noderef{StableColoring}\), and fixes every exact identity incident
with \(U\) in the sense of \(\noderef{TreeIncidentEdge}\).             \(\tag{5}\)

We next verify, rather than merely assert, that every old walk \(Q_i\)
is an exact \((U,\mu_3,\{\gamma,\beta\})\)-exit path.  Consider an
edge-node \(f\) on \(Q_i\).  If \(f\) is not incident with \(U\), the
definition in (4) changes an old \(\alpha\)-value to \(\gamma\), fixes
an old \(\beta\)-value by (3) and the counterexample equality
\(\alpha\ne\beta\), and preserves the relevant incidence adjacency in
\(\noderef{TwoColorIncidenceGraph}\).  Suppose instead that \(f\) is
incident with \(U\).  Since the endpoints of \(Q_i\) are vertex-nodes,
the edge-node \(f\) is internal to the incidence walk and its predecessor
and successor are its two endpoint vertices.  They are distinct by
\(\noderef{LooplessMultigraph}\) and the endpoint-incidence definition
\(\noderef{Incident}\).  The unique-contact clause of
\(\noderef{TreeExitPath}\) says that the only one of these vertices that
can lie in \(V(U)\) is \(b_i\).  Hence exactly one endpoint lies in
\(V(U)\), so
\[
                       f\in\noderef{BoundaryEdges}(G,V(U)).   \tag{6}
\]
Closedness, together with the tree-missingness of \(\alpha\) and
\(\gamma\), excludes both of those old colors on (6).  Since \(f\) is
an old \(\alpha,\beta\)-walk edge, it is therefore \(\beta\)-colored,
and (5) fixes it.

The outside-end missing-color clause also transports field by field.
If \(a_i\) was \(\beta\)-missing, no edge acquires or loses color
\(\beta\), so it remains so.  Suppose it was \(\alpha\)-missing and,
for a contradiction, that an edge \(g\) incident with \(a_i\) has new
color \(\gamma\).  If \(g\) is not incident with \(U\), (4) says that
it had old color \(\alpha\), contradicting the old missingness.  If it
is incident with \(U\), it is fixed.  Since \(a_i\notin V(U)\), its
other incident witness lies in \(V(U)\); the loopless endpoint law makes
\(g\) a genuine boundary identity.  Its fixed old \(\gamma\)-color then
contradicts closedness and the \(\gamma\)-missing witness \(b_3\).
Thus \(a_i\) is \(\gamma\)-missing under \(\mu_3\).  Simplicity,
length, endpoints, support, and unique contact never changed.  We have
proved
\[
 Q_i:\noderef{TreeExitPath}(U,\mu_3,\gamma,\beta,a_i,b_i,Q_i)
                    \qquad(i=1,2,3),                         \tag{7}
\]
with the same pairwise-disjoint supports.  The exact identities
\(f_i\) remain \(\beta\)-colored boundary edges of \(U\), because they
are incident with \(U\) and hence fixed by (5).

We now obtain the first actual replacement ETT.  Restricting (5) along
the exact prefix in (1) gives
\[
                 \mu_3\text{ is }(T_n,D_n,\sigma)\text{-stable}. \tag{8}
\]
For the edge clause, an edge incident with \(T_n\) is incident with
\(U\); for the vertex clause, use the vertex-set inclusion from (1).
The assumed current assertion (vi), applied to \(\mathcal R\), its
maximum property \(\mathcal P.\mathit{maximum}\), and (8), gives
\[
                        M_3:\noderef{ModColoring}(\mathcal R,\mu_3).
                                                                    \tag{9}
\]

The exact terminal route of \(\mathcal P\) also transports from
\(\sigma\) to \(\mu_3\).  At each step of the
\(\sigma\)-\(\noderef{TashkinovAugmentationSequence}\), the added-edge,
missing-color-witness, and exact-prefix clauses of
\(\noderef{TashkinovAugmentation}\) say that the added edge is an edge of
\(U\) incident with a preceding displayed vertex and is therefore fixed
by (5).  The missing-color witness belongs to that preceding tree, hence
to \(V(U)\), where (5) preserves the entire missing set.  Induction over
the reflexive-transitive augmentation chain reuses
the identical ordered steps under \(\mu_3\).  The same vertexwise and
edgewise equalities show that the terminal closedness in (1) is preserved.
Thus the PE closure from \(T_n\), or the SE/RE closure from its retained
exact one-edge seed, is an identical route under \(\mu_3\).

This route is a strict extension of \(T_n\): the vertex \(b_2\) belongs
to \(V(U)\) by (7) but not to \(V(T_n)\) by hypothesis, whereas (1) is
an exact prefix.  The universal clause of (9) therefore yields
\[
 \mathcal A_3:
 \noderef{CorrespondingExtendedTashkinovTree}
       (\mathcal R,\mu_3,U).                                  \tag{10}
\]
By \(\noderef{MaximumPropertyTransport}\), applied to the same
mod-coloring (9), the ETT
\(\mathcal R_3=\mathcal A_3.\mathit{replacementETT}\) has maximum
property.  Apply
\(\noderef{CorrespondingExtendedTashkinovTreeTransTrees}\) to
\(\mathcal P.\mathit{corresponding}\) and \(\mathcal A_3\).  This
gives an actual corresponding certificate from the fixed original
\(\mathcal E\) directly to the \(\mu_3\)-presentation on \(U\).  The
cited theorem composes, in particular, every positive-stage stable
coloring relation through the intermediate replacement; no provenance is
being inferred from equality of the displayed trees.                         \(\tag{11}\)

We identify the exact component to be switched.  Put
\[
 K_3=\noderef{TwoColorComponent}(G,\mu_3,\gamma,\beta,
                                      \operatorname{inl}b_3).
\]
The missingness in (2) is preserved at \(b_3\) by (5), while \(f_3\)
is an incident \(\beta\)-edge.  Properness allows no second incident
\(\beta\)-edge, so \(b_3\) has degree exactly one in the
\(\gamma,\beta\)-incidence graph.  At \(a_3\), (7) gives a missing
color and the nontrivial path gives an incident edge of the other color;
properness again makes its incidence degree exactly one.  Every internal
node of the simple walk \(Q_3\) has its two distinct walk neighbors, and
\(\noderef{TwoColorIncidenceDegreeAtMostTwo}\) excludes a third.
The spanning path in
\(\noderef{TwoColorComponentSpanningPath}\), oriented from the
degree-one node \(a_3\), is therefore forced one node at a time to agree
with \(Q_3\); neither endpoint has room for a continuation.  Hence
\[
                              K_3=Q_3.\mathit{support}.        \tag{12}
\]
Equivalently, a shortest route from a hypothetical node of
\(K_3\setminus Q_3.\mathit{support}\) to the displayed support would
give a third neighbor at an internal node or a second neighbor at one of
the two degree-one endpoints.  This also makes clear that (12) is proved
after the recoloring; it is not imported from the old coloring.

Let \(s\) be the unique displayed index of \(b_3\) in \(U\), whose
existence and uniqueness follow from
\(\noderef{TreeSequenceVertexSet}\) and the vertex-injectivity field of
\(\noderef{TreeSequence}\).  Expanding
\(\noderef{TreeVertexOrder}\), let \(t_2<s\) be the index of \(b_2\).
If \(\ell=|T_n|\), exact prefixhood in (1) and
\(b_2\notin V(T_n)\) imply \(\ell<t_2<s\).  In particular \(s\ge2\).
Define
\[
 B^-=\noderef{TreeSequencePrefix}(U,s-1),\qquad
 C=\noderef{TreeSequencePrefix}(U,s).                         \tag{13}
\]
Then
\[
              T_n\preccurlyeq B^-\prec C\preccurlyeq U,      \tag{14}
\]
the last vertex of \(C\) is \(b_3\), and \(|C|=|B^-|+1\).
Moreover, in the SE/RE route the retained one-edge seed has length
\(\ell+1\le s-1=|B^-|\), so that seed too is an exact prefix of \(B^-\).

By (7), \(b_3\) is the only vertex-node of \(Q_3\) in \(V(U)\).
Equations (12)--(14) therefore show that \(K_3\) contains no vertex-node
of \(B^-\).  If an edge-node of \(K_3\) were incident with \(B^-\),
the incidence adjacency would put its endpoint vertex-node in the same
component, a contradiction.  Thus no edge identity in \(K_3\) is
incident with \(B^-\).  Define
\[
 \mu_4=\noderef{TwoColorComponentSwap}
          (G,\mu_3,\gamma,\beta,\operatorname{inl}b_3).       \tag{15}
\]
By (3) and \(\noderef{TwoColorComponentSwapProper}\), \(\mu_4\) is
proper.  The swap transposes only the two displayed nonempty values, so
the two clauses of \(\noderef{IsOneEdgeUncolored}\) show that precisely
\(e_0\) remains uncolored.  Every edge incident with \(B^-\) is fixed;
using \(\noderef{PresentColors}\) and \(\noderef{MissingColors}\),
every missing set at a vertex of \(B^-\) is fixed as well.  In particular
\[
                  \mu_4\text{ is }(T_n,D_n,\mu_3)\text{-stable}. \tag{16}
\]
The correspondence fields in (10) identify its terminal tree with
\(T_n\), its terminal series coloring with \(\mu_3\), and every raw
stored set with the corresponding raw set of \(\mathcal R\).  We also
check the computed terminal stored set.  By
\(\noderef{TashkinovStoredColors}\), it is the union of the raw stored
sets through the terminal index, with the colors missing on the terminal
ladder deleted.  The raw unions agree by the stored-set fields of (10).
The ladder trees agree, and the vertex clause of the positive-stage
stability field in (10) identifies the two
\(\noderef{MissingColors}\) sets at every vertex of that ladder; taking
their union as in \(\noderef{TreeMissingColors}\) identifies the deleted
sets.  Hence the computed set for \(\mathcal R_3\) is exactly \(D_n\).
Since \(\mathcal R_3\) has the same rung as \(\mathcal E\) and maximum
property, current assertion (vi) applied to (16) gives
\[
                        M_4:\noderef{ModColoring}(\mathcal R_3,\mu_4).
                                                                    \tag{17}
\]

We reconstruct the exact route through \(b_3\).  The transported
\(\mu_3\)-route used in (10), restricted at the exact prefix \(C\), is
a route from \(T_n\) in PE, or from the retained exact one-edge seed in
SE/RE, to \(C\); the restriction is precisely
\(\noderef{TashkinovAugmentationSequenceRestrict}\), using (14).
The component swap fixes every displayed edge of \(C\): an earlier edge
is incident with \(B^-\), and the final displayed edge attaching \(b_3\)
also has its old endpoint in \(B^-\).  It also fixes every missing set on
\(B^-\).  Apply
\(\noderef{TashkinovAugmentationSequenceColoringTransport}\) with target
\(C\) and penultimate prefix \(B^-\).  Its two pointwise premises are
exactly these fixed-edge and fixed-missing-set facts, and it yields the
identical route under \(\mu_4\).                                  \(\tag{18}\)

Choose a finite closure
\[
                    \noderef{TashkinovClosure}(\mu_4,C,T')    \tag{19}
\]
by \(\noderef{TashkinovClosureExists}\).  Concatenating its augmentation
sequence with (18) gives the exact PE or SE/RE terminal route to \(T'\).
It is strict over the terminal ladder because \(C\) already contains the
vertex \(b_2\notin V(T_n)\).  The universal clause of (17) therefore
supplies
\[
 \mathcal A_4:
 \noderef{CorrespondingExtendedTashkinovTree}
       (\mathcal R_3,\mu_4,T').                               \tag{20}
\]
Apply \(\noderef{MaximumPropertyTransport}\) to \(M_4\), (20), and the
maximum property of \(\mathcal R_3\); its replacement ETT has maximum
property under \(\mu_4\).                                    \(\tag{21}\)

There are two required provenance compositions.  First apply
\(\noderef{CorrespondingExtendedTashkinovTreeTransTrees}\) to (10) and
(20), producing a certificate from \(\mathcal R\) to \(\mu_4,T'\).
Then apply it again to \(\mathcal P.\mathit{corresponding}\) and that
certificate.  The result is an actual certificate
\[
 \mathcal A_{40}:
 \noderef{CorrespondingExtendedTashkinovTree}
       (\mathcal E,\mu_4,T')                                  \tag{22}
\]
relative to the fixed original ETT.  At each application the theorem
returns heterogeneous equality with the final nested ETT and series;
transporting (21) along that equality gives maximum property for the
replacement ETT in (22).  Likewise, its ladder-tree, kind, and exact
connecting-edge equalities transport the concatenated route preceding
(20) to the terminal data of (22).  Thus (19), (21), and (22), with the
same PE route or the same retained SE/RE edge identity and seed, form a
\(\noderef{RutgersClosedPresentation}\) \(\mathcal P'\) relative to
\(\mathcal E\).  This use of maximum property comes from the applicable
mod-coloring (17); neither maximum property nor correspondence provenance
is inferred merely from equality of displayed trees.

We now build the shorter counterexample on \(\mathcal P'\).  The swap in
(15) preserves the entire underlying \(\gamma,\beta\)-incidence graph:
inside \(K_3\) it interchanges the two allowed colors, and outside
\(K_3\) it fixes every edge.  Therefore the same three node lists are
walks in that graph.  For \(Q_3\), the missing statuses at its two
degree-one endpoints are interchanged.  The other two walk supports lie
outside \(K_3\): if their components met \(K_3\), connected components
would coincide, contrary to (12) and the retained support-disjointness;
this is the disjointness principle in
\(\noderef{TwoColorComponentsDisjoint}\).  Their edges and endpoint
missing statuses are fixed.  Hence all three statements in (7) remain
exact \(\gamma,\beta\)-exit statements under \(\mu_4\).  At \(b_3\),
the unique incident edge of \(K_3\), namely \(f_3\), changes from
\(\beta\) to \(\gamma\), while no old incident \(\gamma\)-edge exists by
(2) and (5).  Consequently
\[
                 \beta\in\noderef{MissingColors}(G,\mu_4,b_3).
                                                                    \tag{23}
\]

The prefix \(C\), hence the closure \(T'\), contains all three \(b_i\).
The identities \(f_1,f_2\) are fixed by (15), because their edge-nodes
lie in the supports of \(Q_1,Q_2\), and remain \(\beta\)-colored.  We
claim that both endpoints of each belong to \(V(T')\).  One endpoint
\(b_i\) already belongs to \(V(C)\).  If its other endpoint were outside
\(V(T')\), the loopless endpoint law and
\(\noderef{BoundaryEdges}\) would make \(f_i\) a \(\beta\)-colored
boundary identity of \(T'\).  But (23) makes \(\beta\) tree-missing on
\(T'\), and the closedness clause of (19) forbids such an identity.
This proves the claim, in particular for both ends of the retained exact
identity \(f_1\).                                               \(\tag{24}\)

Apply the assumed current assertion (i) to the replacement ETT in (22),
using its rung agreement and transported maximum property.  It gives
\[
       \noderef{ColoredVertexSetElementary}(G,\mu_4,V(T')).   \tag{25}
\]
No outside end \(a_i\) lies in \(V(T')\).  Suppose one did.  It is
distinct from \(b_3\), since it was outside \(U\) while \(b_3\in V(U)\).
One of \(\gamma,\beta\) is missing there by the transported exit
certificate, whereas \(\beta\) is missing at \(b_3\) by (23).
Elementariness (25) therefore forces \(\gamma\) to be missing at that
\(a_i\).  Both \(\gamma\) and \(\beta\) are now tree-missing on the
closed set \(V(T')\).  Traverse each \(Q_j\) from
\(b_j\in V(C)\subseteq V(T')\) toward \(a_j\).  If it ever first left
\(V(T')\), the crossing edge would be a \(\gamma\)- or
\(\beta\)-colored boundary identity, forbidden by closedness.  Thus all
three distinct \(a_j\) lie in \(V(T')\).  At none of them can \(\beta\)
be missing, by (23), (25), and their distinctness from \(b_3\); hence
\(\gamma\) is missing at all three.  This contradicts (25) already at
\(a_1,a_2\).  Therefore
\[
                              a_i\notin V(T')\quad(i=1,2,3).   \tag{26}
\]

For each \(i\), traverse the finite simple path \(Q_i\) from \(a_i\)
and stop at its first vertex-node \(c_i\) in \(V(T')\).  Such a node
exists because \(b_i\in V(T')\); (26) says it is not the initial node.
Let \(L_i\) be this initial subwalk and let \(g_i\) be its final exact
edge identity.  The predecessor vertex of \(g_i\) lies outside
\(V(T')\), while \(c_i\) lies inside.  The two are its distinct endpoints
by \(\noderef{Incident}\) and \(\noderef{LooplessMultigraph}\), so
\[
                         g_i\in\noderef{BoundaryEdges}(G,V(T')). \tag{27}
\]
Its incidence adjacency says that its color is \(\gamma\) or \(\beta\).
The second alternative is forbidden by closedness and (23); hence
\[
                         \mu_4(g_i)=\gamma.                   \tag{28}
\]

Equations (26)--(28) verify every field of
\(\noderef{TreeExitPath}(T',\mu_4,\beta,\gamma,a_i,c_i,L_i)\):
the color order is immaterial to the incidence graph; the subwalk is
simple and nontrivial, starts outside, ends at its unique first tree
vertex, and retains its endpoint missing-color alternative.  Its support
is contained in that of \(Q_i\), so the three supports remain pairwise
disjoint and their outside ends remain pairwise distinct.  The new tree
ends \(c_i\) are pairwise distinct.  If \(c_i=c_j\), then support
disjointness gives \(g_i\ne g_j\); both identities are incident with the
common vertex and hence adjacent by \(\noderef{AdjacentEdges}\), while
(28) gives them the same nonempty color, contradicting properness.

Equation (23) proves
\(\beta\in\noderef{TreeMissingColors}(\mu_4,T')\).  On the other hand
\(\gamma\notin\noderef{TreeMissingColors}(\mu_4,T')\): if \(\gamma\)
were missing somewhere on \(T'\), closedness would forbid the exact
\(\gamma\)-colored boundary identity \(g_1\) from (27)--(28).
Thus the formal normalization required by
\(\noderef{RutgersCounterexample}\) is
\[
                \alpha'=\beta\quad\hbox{(missing)},\qquad
                \beta'=\gamma\quad\hbox{(not tree-missing)}. \tag{29}
\]
This is the same unordered two-color pair used in the paper, oriented so
that the structure's first color is its missing color.  Distinctness is
(3), and (27)--(28) give the required exact \(\beta'\)-colored boundary
identities \(g_i\), incident with \(c_i\) and lying on \(L_i\).

Use \(\noderef{TreeSequenceVertexSet}\) and vertex injectivity in
\(\noderef{TreeSequence}\) to obtain the three unique displayed indices
of the distinct \(c_i\).  Sort the complete packages
\((a_i,c_i,L_i,g_i)\) by those indices.  The strict linear index order is
exactly \(\noderef{TreeVertexOrder}\), so after relabeling the tree ends
satisfy \(c_1\prec_{T'}c_2\prec_{T'}c_3\).  Permuting whole packages
preserves all support, distinctness, exit, incidence, and boundary fields.
Together with (22), (25), and (29), these data define a Rutgers
counterexample \(\mathcal C'\) on \(\mathcal P'\).

Finally, each \(L_i\) is an initial subpath of \(Q_i\), so
\(|L_i|\le |Q_i|\).  The inequality is strict for the package containing
the old \(Q_1\).  Indeed, (24) places both endpoint vertices of its
retained edge-node \(f_1\) in \(V(T')\), while \(f_1\) belongs to
\(Q_1.\mathit{support}\).  A path traversed from \(a_1\notin V(T')\)
must therefore meet an endpoint vertex in \(V(T')\) before it could
include the intervening edge-node \(f_1\).  The first-contact subpath
stops at that vertex and omits \(f_1\); hence \(|L_1|<|Q_1|\).
Sorting does not change the sum, so the definition
\(\noderef{RutgersCounterexampleMeasure}\) gives
\[
 \noderef{RutgersCounterexampleMeasure}(\mathcal C')
   =|L_1|+|L_2|+|L_3|
   <|Q_1|+|Q_2|+|Q_3|
   =\noderef{RutgersCounterexampleMeasure}(\mathcal C).
\]
This proves the asserted strict shortening without any appeal to
minimality.
\end{proof}
